\chapter{Review and Clean Up}

We began with an empty topology and finished with a testable system of five
jails. Each boundary has a distinct responsibility, and Blackship knows the
dependency order between them:

\begin{lstlisting}
test -> web -> store -> catalog -> database
\end{lstlisting}

The web jail is the only application component with a published inbound port.
The Rails applications and PostgreSQL communicate over the private BSDRunner
network, while the test jail behaves as an external client of the private web
edge.

\section{What the final test proves}

One invocation of \texttt{rails-on-jails-test} now performs a compact but real
distributed-system check:

\begin{enumerate}
  \item Blackship starts any missing jails in dependency order.
  \item The readiness loop waits until nginx can return a successful page.
  \item The test jail creates a unique product through catalog's HTTP API.
  \item Catalog persists that product in PostgreSQL.
  \item A request enters through nginx and reaches the store Rails app.
  \item Store calls catalog over HTTP and renders the returned product.
  \item The test observes its unique name, SKU, and formatted price in HTML.
  \item The fixture is deleted through the catalog API.
  \item The host receives status zero on success and nonzero when the deployed
    graph cannot become ready.
\end{enumerate}

This is stronger than a collection of isolated unit tests because the result
depends on live processes, routing, configuration, serialization, persistence,
and presentation all agreeing. It does not replace unit tests: those remain
the faster place to cover application edge cases.

\section{What remains outside the test}

The experiment intentionally does not claim to test TLS, authentication,
JavaScript in a browser, concurrent writers, backup restoration, rolling
deployments, or high availability. It also uses fixed private addresses and a
single FreeBSD host. Those are reasonable future exercises, not hidden
properties of the test we built.

The self-owned fixture makes repeated runs independent, but an abrupt process
termination between creation and teardown could still leave a product behind.
A production suite can add a run identifier and a separate stale-fixture
cleanup policy without weakening the normal teardown path.

\section{Exercise 27A: Record the final checkpoint}

With every service restored, run the host command one final time and capture
its status in an explicit shell:

\begin{lstlisting}[language=sh]
sh -c '
  mdo /usr/local/sbin/rails-on-jails-test
  code=$?
  echo "exit status: $code"
  exit "$code"
'
\end{lstlisting}

The completed checkpoint is one run, nine assertions, no failures or errors,
and \texttt{exit status: 0}. Keep that output with any notes about the host's
FreeBSD and Blackship versions; generated credentials do not belong in the
record.

\section{Exercise 27B: Choose a final state}

No cleanup is required to finish the tutorial. Choose one of these states
deliberately:

\begin{description}
  \item[Keep running] Leave the graph available for more experiments.
  \item[Stop and preserve] Stop \texttt{database} through BSDRunner and accept
    its dependency-aware plan to stop the higher-level jails first. The jail
    filesystems and manifests remain available for a later start.
  \item[Remove permanently] First remove the web jail's host port forward.
    Then delete the jails through BSDRunner in this order: \texttt{test},
    \texttt{web}, \texttt{store}, \texttt{catalog}, and \texttt{database}.
    Review every deletion target before confirming it; deleting database also
    deletes the PostgreSQL data created by this experiment.
\end{description}

If the graph is removed permanently, also remove the host runner:

\begin{lstlisting}
/usr/local/sbin/rails-on-jails-test
\end{lstlisting}

Destroy the private temporary record that held the PostgreSQL password as
well. Do this only after deciding that the database will not be reused. The
tutorial source and generated PDF contain no credential and can be retained.

\section{Where to go next}

Useful extensions include placing TLS at the web edge, authenticating catalog
writes, assigning a per-run namespace for parallel tests, collecting service
logs as failure artifacts, or invoking the host runner from a CI agent that
owns this FreeBSD environment. None requires changing the central testing
idea: arrange the real boundaries, create the test's data through a supported
interface, observe behavior at the edge, and return a truthful process status.

\section{Final checkpoint}

The tutorial is complete when the final host invocation returns zero and a
deliberately unavailable catalog makes the same workflow return nonzero. At
that point we have not merely run Rails in jails; we have built an inspectable,
failure-aware end-to-end test across Ruby, Rails, PostgreSQL, nginx, Blackship,
and FreeBSD.
